% ========================================================
% BÀI TẬP TOÁN 5 - KẾT NỐI TRI THỨC VỚI CUỘC SỐNG
% Bài 9: Luyện tập chung
% Định dạng: BTTeX Part III - Trắc nghiệm điền số ngắn
% ========================================================

% Câu 1
\begin{ex}
  Tính giá trị biểu thức: $45\,000 + 12\,500 \times 4$.
  \shortans{95000}
  \loigiai{
    $12\,500 \times 4 = 50\,000$. Kết quả là $45\,000 + 50\,000 = 95\,000$.
  }
\end{ex}

% Câu 2
\begin{ex}
  Tính giá trị biểu thức: $120\,000 : 60 - 150 \times 4$.
  \shortans{1400}
  \loigiai{
    $120\,000 : 60 = 2\,000$; $150 \times 4 = 600$. Kết quả: $2\,000 - 600 = 1\,400$.
  }
\end{ex}

% Câu 3
\begin{ex}
  Rút gọn phân số $\frac{72}{96}$ thành phân số tối giản. Tổng của tử số và mẫu số là bao nhiêu?
  \shortans{7}
  \loigiai{
    Chia cả tử và mẫu cho 24: $\frac{72 : 24}{96 : 24} = \frac{3}{4}$. Tổng tử số và mẫu số là $3 + 4 = 7$.
  }
\end{ex}

% Câu 4
\begin{ex}
  Chuyển phân số $\frac{7}{4}$ thành phân số thập phân có mẫu số là $100$. Tử số của phân số đó là bao nhiêu?
  \shortans{175}
  \loigiai{
    $\frac{7 \times 25}{4 \times 25} = \frac{175}{100}$. Tử số là 175.
  }
\end{ex}

% Câu 5
\begin{ex}
  Tính giá trị biểu thức: $\frac{5}{6} \times \frac{3}{10} + \frac{3}{4}$. Tử số của kết quả tối giản là bao nhiêu?
  \shortans{1}
  \loigiai{
    $\frac{5 \times 3}{6 \times 10} = \frac{1}{4}$. $\frac{1}{4} + \frac{3}{4} = \frac{4}{4} = 1$. Tử số khi viết là 1 (hoặc kết quả là 1).
  }
\end{ex}

% Câu 6
\begin{ex}
  Tính kết quả: $2\frac{3}{4} + 1\frac{1}{4}$. Kết quả là bao nhiêu?
  \shortans{4}
  \loigiai{
    $(2 + 1) + (\frac{3}{4} + \frac{1}{4}) = 3 + 1 = 4$.
  }
\end{ex}

% Câu 7
\begin{ex}
  Tìm $x$ biết: $x \times 1\frac{1}{2} = 6$. Giá trị của $x$ là bao nhiêu?
  \shortans{4}
  \loigiai{
    $x = 6 : 1\frac{1}{2} = 6 : \frac{3}{2} = 4$.
  }
\end{ex}

% Câu 8
\begin{ex}
  Đổi $5\text{ m}^2\text{ } 8\text{ dm}^2$ thành bao nhiêu đề-xi-mét vuông?
  \shortans{508}
  \loigiai{
    $5\text{ m}^2 = 500\text{ dm}^2$. Do đó $5\text{ m}^2\text{ } 8\text{ dm}^2 = 508\text{ dm}^2$.
  }
\end{ex}

% Câu 9
\begin{ex}
  Đổi $3\text{ tấn } 5\text{ tạ}$ thành bao nhiêu yến?
  \shortans{350}
  \loigiai{
    $3\text{ tấn} = 300\text{ yến}$, $5\text{ tạ} = 50\text{ yến}$. Do đó $3\text{ tấn } 5\text{ tạ} = 350\text{ yến}$.
  }
\end{ex}

% Câu 10
\begin{ex}
  Một sân trường hình chữ nhật có chu vi $200\text{ m}$, chiều rộng kém chiều dài $20\text{ m}$. Diện tích sân trường là bao nhiêu mét vuông?
  \shortans{2400}
  \loigiai{
    Nửa chu vi là $200 : 2 = 100\text{ m}$. Chiều dài là $(100 + 20) : 2 = 60\text{ m}$. Chiều rộng là $100 - 60 = 40\text{ m}$. Diện tích là $60 \times 40 = 2\,400\text{ m}^2$.
  }
\end{ex}

% Câu 11
\begin{ex}
  Một ô tô chở được $50$ bao gạo, mỗi bao nặng $50\text{ kg}$. Xe ô tô đó chở được tất cả bao nhiêu tạ gạo?
  \shortans{25}
  \loigiai{
    Tổng khối lượng là: $50 \times 50 = 2\,500\text{ kg} = 25\text{ tạ}$.
  }
\end{ex}

% Câu 12
\begin{ex}
  Tính nhanh: $37 \times 45 + 37 \times 55$.
  \shortans{3700}
  \loigiai{
    $37 \times (45 + 55) = 37 \times 100 = 3\,700$.
  }
\end{ex}

% Câu 13
\begin{ex}
  Tìm số tự nhiên lớn nhất có bốn chữ số khác nhau chia hết cho cả $2$ và $5$.
  \shortans{9870}
  \loigiai{
    Số tận cùng là 0 để chia hết cho 2 và 5. Chọn các chữ số hàng nghìn, trăm, chục lớn nhất khác nhau là 9, 8, 7. Số đó là $9\,870$.
  }
\end{ex}

% Câu 14
\begin{ex}
  Một mảnh vườn trồng cam và bưởi có tất cả $120$ cây. Số cây cam chiếm $\frac{2}{3}$ tổng số cây. Hỏi mảnh vườn có bao nhiêu cây bưởi?
  \shortans{40}
  \loigiai{
    Số cây bưởi chiếm $1 - \frac{2}{3} = \frac{1}{3}$ số cây. Số cây bưởi là $120 \times \frac{1}{3} = 40$ cây.
  }
\end{ex}

% Câu 15
\begin{ex}
  Tìm một số biết rằng nếu lấy số đó chia cho $4$ rồi cộng với $35$ thì được $60$. Số đó là bao nhiêu?
  \shortans{100}
  \loigiai{
    Gọi số đó là $x$. Ta có $x : 4 + 35 = 60 \implies x : 4 = 25 \implies x = 100$.
  }
\end{ex}

% Câu 16
\begin{ex}
  Trung bình cộng của ba số là $45$. Biết số thứ nhất là $38$, số thứ hai là $42$. Số thứ ba là bao nhiêu?
  \shortans{55}
  \loigiai{
    Tổng của ba số là $45 \times 3 = 135$. Số thứ ba là $135 - 38 - 42 = 55$.
  }
\end{ex}

% Câu 17
\begin{ex}
  Tính giá trị biểu thức: $125 \times 8 \times 25 \times 4$. Kết quả là bao nhiêu?
  \shortans{100000}
  \loigiai{
    $(125 \times 8) \times (25 \times 4) = 1\,000 \times 100 = 100\,000$.
  }
\end{ex}

% Câu 18
\begin{ex}
  Cho phân số $\frac{21}{35}$. Rút gọn về phân số tối giản rồi nhân với $15$. Kết quả là bao nhiêu?
  \shortans{9}
  \loigiai{
    $\frac{21}{35} = \frac{3}{5}$. Lấy $\frac{3}{5} \times 15 = 9$.
  }
\end{ex}

% Câu 19
\begin{ex}
  Một hình thoi có diện tích $48\text{ cm}^2$, độ dài một đường chéo là $12\text{ cm}$. Độ dài đường chéo còn lại là bao nhiêu xăng-ti-mét?
  \shortans{8}
  \loigiai{
    Độ dài đường chéo còn lại là: $(48 \times 2) : 12 = 8\text{ cm}$.
  }
\end{ex}

% Câu 20
\begin{ex}
  Một người thợ làm việc từ $7\text{ giờ } 30\text{ phút}$ đến $11\text{ giờ } 30\text{ phút}$ và làm được $8$ sản phẩm. Hỏi trung bình mỗi sản phẩm người đó làm trong bao nhiêu phút?
  \shortans{30}
  \loigiai{
    Thời gian làm là $11\text{ giờ } 30\text{ phút} - 7\text{ giờ } 30\text{ phút} = 4\text{ giờ} = 240\text{ phút}$. Trung bình mỗi sản phẩm làm trong $240 : 8 = 30$ phút.
  }
\end{ex}

% Câu 21
\begin{ex}
  Tìm số tự nhiên bé nhất có năm chữ số khác nhau chia hết cho cả $2, 5$ và $9$.
  \shortans{12690}
  \loigiai{
    Số chia hết cho 2 và 5 nên chữ số tận cùng là 0. Để số chia hết cho 9 thì tổng các chữ số chia hết cho 9. Để số là bé nhất, ta chọn hàng chục nghìn là 1, hàng nghìn là 2, hàng trăm là 6, hàng chục là 9 (vì $1 + 2 + 6 + 9 + 0 = 18$). Vậy số bé nhất là $12\,690$.
  }
\end{ex}

% Câu 22
\begin{ex}
  Có bao nhiêu số chẵn có hai chữ số chia hết cho $3$?
  \shortans{15}
  \loigiai{
    Số chẵn chia hết cho 3 là số chia hết cho 6. Các số có hai chữ số chia hết cho 6 là: $12, 18, 24, \dots, 96$. Số lượng số là: $(96 - 12) : 6 + 1 = 15$ số.
  }
\end{ex}

% Câu 23
\begin{ex}
  Tính tổng của tất cả các phân số có mẫu số bằng $5$ và nhỏ hơn $1$ (tử số lớn hơn $0$).
  \shortans{2}
  \loigiai{
    Các phân số là: $\frac{1}{5} + \frac{2}{5} + \frac{3}{5} + \frac{4}{5} = \frac{10}{5} = 2$.
  }
\end{ex}

% Câu 24
\begin{ex}
  Một hình chữ nhật có chu vi $72\text{ cm}$. Nếu tăng chiều rộng thêm $6\text{ cm}$ và giữ nguyên chiều dài thì được một hình vuông. Diện tích hình chữ nhật ban đầu là bao nhiêu xăng-ti-mét vuông?
  \shortans{315}
  \loigiai{
    Nửa chu vi là $72 : 2 = 36\text{ cm}$. Chiều dài hơn chiều rộng là $6\text{ cm}$. Chiều dài là: $(36 + 6) : 2 = 21\text{ cm}$. Chiều rộng là: $36 - 21 = 15\text{ cm}$. Diện tích ban đầu là: $21 \times 15 = 315\text{ cm}^2$.
  }
\end{ex}

% Câu 25
\begin{ex}
  Một đàn vịt có một số con đang bơi dưới ao và một số con trên bờ. Số vịt dưới ao gấp 3 lần số vịt trên bờ. Nếu có 3 con từ dưới ao lên bờ thì số vịt dưới ao chỉ còn gấp đôi số vịt trên bờ. Hỏi cả đàn có bao nhiêu con vịt?
  \shortans{36}
  \loigiai{
    Ban đầu số vịt trên bờ bằng $\frac{1}{4}$ cả đàn. Khi 3 con lên bờ, số vịt trên bờ bằng $\frac{1}{3}$ cả đàn. 3 con vịt ứng với: $\frac{1}{3} - \frac{1}{4} = \frac{1}{12}$ cả đàn. Cả đàn vịt có: $3 : \frac{1}{12} = 36$ (con).
  }
\end{ex}
